% ============================================================
% THEO-CHIR-PCD-ORIENTATION-1: The n-hat -> omega_PCD Link as a
%   Chirality-Classification Theorem
% A Viability-Level Primitive-Count Theorem (resolves audit E20 /
%   OPEN-CHIR-F1-LINK)
% Conscious Point Physics --- Substrate Chirality Arc
%   (series_umbrella/series_substrate_chirality_arc/chirality_derivations/)
%
% Purpose: resolve the classification of audit entry E20 (the
%   n-hat -> omega_PCD link, audit dynamics sub-question D4), left at
%   "unregistered (conditional)" by THEO-CHIR-AUDIT-1 v1.1. The theorem
%   answers the PRIMITIVE-COUNT question (not the magnitude question):
%   does the PCD-cycle orientation pseudovector omega_PCD introduce an
%   independent chirality primitive, or does it reduce to the chirality
%   primitives the audit already registered? Finding: omega_PCD =
%   sigma_cycle * n-hat is a product of the registered spatial primitive
%   n-hat (E16) and the registered temporal primitive's handedness
%   sigma_cycle (E2/E5/E17, carried by A1+A4). No independent third
%   primitive: Scenario B refuted; E20 reclassified emergent (provisional).
%
% Scope/layer: VIABILITY LEVEL (Layer 2.5). The theorem lifts the F.1
%   sub-question sketch's Phase-1 (net DI-bit current || n-hat) and Phase-2
%   (coupling rule) results, which are themselves viability-level with three
%   open structural commitments (F.1 sketch section 14.17). The theorem does
%   NOT exceed that ceiling: it caps E20 at emergent (provisional), not
%   established. It does NOT close the full F.1 Scenario-A trajectory, does
%   NOT derive the magnitude |chi|=phi^-3 (that is OPEN-CHIR-1d / E21), and
%   does NOT derive n-hat or sigma_cycle (both remain primitives).
%
% Patch 0636 (Session 148) --- first downstream derivation theorem of the
%   audit-spawned OPEN-CHIR-* programme. Built from the Patch 0635 scope
%   sketch theo_chir_pcd_orientation_1_scope.md. Precondition cleared
%   (section 3): the F.1 sketch's "sigma_cycle from axiom A5" attribution is a
%   pre-canonical-numbering drift (canonical A5 = metric eta=1/phi); the
%   cycle ordering is A1+A4 content (the audit temporal primitive). This
%   artifact attributes sigma_cycle to A1+A4.
%
% CHANGELOG
%   v1.0 (Patch 0636, Session 148, 29 May 2026) --- initial registration of
%     THEO-CHIR-PCD-ORIENTATION-1 at viability-level (Layer 2.5) as a
%     primitive-count classification theorem resolving audit E20. Scenario B
%     refuted; E20 -> emergent (provisional). Six sections per scope-sketch
%     section 4.
% ============================================================
% Version 1.1 --- 17 August 2026 (Patch 3212):
%   added the generated plain-language appendix glossing the internal
%   programme identifiers used in this paper (founder jargon ruling, 16
%   Aug 2026). No claim, equation, number or section changed.

\documentclass[11pt,letterpaper]{article}
\usepackage[utf8]{inputenc}
\usepackage[margin=1in]{geometry}
\usepackage[T1]{fontenc}
\usepackage{amsmath,amssymb,amsfonts,amsthm}
\usepackage{xcolor}
\usepackage{hyperref}
\usepackage{enumitem}
\usepackage{parskip}
\usepackage{mdframed}
\usepackage{longtable}
\usepackage{array}

\hypersetup{
    colorlinks=true,
    linkcolor=blue,
    citecolor=blue,
    urlcolor=blue
}

\newtheorem{theorem}{Theorem}[section]
\newtheorem{proposition}[theorem]{Proposition}
\newtheorem{lemma}[theorem]{Lemma}
\newtheorem{corollary}[theorem]{Corollary}
\newtheorem{definition}[theorem]{Definition}
\newtheorem{remark}[theorem]{Remark}

\newcommand{\nhat}{\hat{n}}
\newcommand{\vhost}{v_{\text{host}}}
\newcommand{\Reals}{\mathbb{R}}
\newcommand{\Hthree}{H_3}
\newcommand{\Hfour}{H_4}
\newcommand{\Icoh}{I_h}
\newcommand{\Oof}[1]{\mathcal{O}(#1)}
\newcommand{\wPCD}{\vec{\omega}_{\text{PCD}}}
\newcommand{\jnet}{\hat{j}^{\text{net}}}
\newcommand{\scyc}{\sigma_{\text{cycle}}}

\title{\textbf{THEO-CHIR-PCD-ORIENTATION-1}\\[4pt]
The $\nhat \mapsto \wPCD$ Link as a Chirality-Classification Theorem\\[6pt]
\large A Viability-Level Primitive-Count Resolution of Audit Entry E20}
\author{Conscious Point Physics --- Substrate Chirality Arc}
\date{Version 1.0 --- Patch 0636, Session 148 --- 29 May 2026\\[2pt]
{\small Version~1.1 --- 17 August 2026 (Patch 3212: added the generated plain-language appendix glossing the internal programme identifiers used in this paper (founder jargon ruling, 16 Aug 2026). No claim, equation, number or section changed.)}}

\begin{document}
\maketitle

\begin{abstract}
\noindent
THEO-CHIR-AUDIT-1 (v1.1) left exactly one entry of its 27-entry chirality
classification map undecided: entry E20, the link between the substrate primitive 4D
direction $\nhat$ and the PCD-cycle orientation pseudovector $\wPCD$. The audit
classified E20 \emph{unregistered (conditional)} --- emergent under Scenario~A
($\wPCD$ derived from $\nhat$) and primitive under Scenario~B ($\wPCD$ an independent
primitive) --- and registered the resolution as OPEN-CHIR-F1-LINK. We resolve the
classification by answering the \emph{primitive-count} question rather than the
magnitude question. Lifting the F.1 sub-question sketch's Phase-1 result (the substrate
net DI-bit current is aligned with $\nhat$) and Phase-2 result (the cycle-orientation
coupling rule), the cycle-orientation pseudovector decomposes as $\wPCD = \scyc\,\nhat$,
the product of the registered spatial primitive $\nhat$ (audit E16, the axis) and the
registered temporal primitive's handedness $\scyc$ (audit E2/E5/E17, the sign). We
reconcile a pre-canonical axiom-attribution drift: $\scyc$ is content of A1$+$A4 (the
temporal primitive), not of A5 (the metric). The decomposition shows $\wPCD$ introduces
\emph{no independent third chirality primitive}: \textbf{Scenario~B is refuted and E20
is emergent from the two registered primitives jointly}, preserving the audit's headline
that all spatial chirality reduces to $\nhat$. The result is stated at \emph{viability
level} (Layer~2.5): it inherits the F.1 trajectory's three open structural commitments
as a ceiling, so E20 reclassifies to \emph{emergent (provisional)}, not established. The
load-bearing economy is that the primitive-count is robust to all three open commitments
(none reintroduces an independent direction or handedness), even though the magnitude
and Layer-3 rigor remain provisional. Falsification is by exhibiting an independent
direction in $\wPCD$ unequal to $\pm\nhat$, or a sign content independent of the temporal
primitive's ordering.
\end{abstract}

\tableofcontents

% ============================================================
\section{Setup: the primitive-count question}
\label{sec:setup}
% ============================================================

\subsection{What E20 is, and what this theorem must decide}

THEO-CHIR-AUDIT-1 catalogued the PCD-cycle orientation pseudovector $\wPCD(\vhost)$ ---
the axial vector representing which way the Perceive $\to$ Compute $\to$ Displace cycle
progresses at each Conscious Point at each Absolute Moment --- as audit entry E20
(dynamics-pass sub-question D4). At v1.1 the audit classified E20 \emph{unregistered
(conditional)}: \emph{emergent} under Scenario~A ($\wPCD$ derived from $\nhat$ via a
substrate mechanism) and \emph{primitive} under Scenario~B ($\wPCD$ an independent
primitive parallel to $\nhat$). The audit did not pre-commit (exclusion X2) and
registered the resolution as OPEN-CHIR-F1-LINK.

The question this theorem decides is the \textbf{primitive-count} question:

\begin{quote}
Does $\wPCD$ introduce an independent chirality primitive --- a third foundational input
parallel to $\nhat$ and the temporal primitive --- or does it reduce to the chirality
primitives the audit has already registered?
\end{quote}

This is deliberately narrower than the full F.1 Scenario-A closure. It is a question about
\emph{how many independent primitives} the framework carries, not about the magnitude of
$\wPCD$, the leptogenesis empirical check, or the Layer-3 derivation of the coupling. As
\S\ref{sec:count} shows, the narrow question admits a more robust answer than the full
closure, and that robustness is the reason E20 is the right first downstream theorem.

\subsection{The two registered primitives the audit already carries}

The audit established two irreducible chirality primitives, against which E20 is to be
measured:

\begin{description}
\item[The spatial primitive $\nhat$ (audit E16, G-6).] The substrate primitive 4D
direction, FI-C-RC-1, established at Capotauro v2.0 Reading C: under vertex-aligned
$\nhat = \vhost$, the $\Hfour$ symmetry breaks to $\Hthree = \Icoh$ at the host vertex.
This is the \emph{unique spatial primitive}; the audit's headline finding is that all
spatial chirality reduces to it.
\item[The temporal primitive (audit E2/E5/E17).] The ordered PCD sequence, carried by
A1 (the CP's capacity to perceive and respond, which defines the cycle's steps) $+$ A4
(the Nexus / Absolute-Moment cadence, which sequences them). Per the audit's v1.1
clarification, the \emph{ordering} of the cycle is primitive; its time-reversal
\emph{asymmetry} (irreversibility) is a logically distinct property registered separately
as OPEN-CHIR-2a. The cycle's handedness --- which of P$\to$C$\to$D or D$\to$C$\to$P is the
forward direction --- is the pseudoscalar content of this temporal primitive.
\end{description}

The theorem will show $\wPCD$ is a function of exactly these two, and no third.

\subsection{Operational senses $\wPCD$ carries}

$\wPCD$ is an axial vector. By the audit's three operational senses, it carries
\emph{spatial} content (its axis is a direction in the substrate's $\Reals^4$) and
\emph{temporal} content (its sign encodes the cycle's forward direction, a $T$-odd datum).
It is not in itself a CP-asymmetric object; CP-asymmetry enters downstream when $\wPCD$
combines with matter content (audit E22/E23/E26). The classification therefore concerns
the spatial axis and the temporal sign separately --- which is exactly the factorization
the decomposition exhibits.

% ============================================================
\section{The decomposition (lifted from the F.1 sub-question)}
\label{sec:decomp}
% ============================================================

This theorem does not re-derive the substrate-mechanism content; the F.1 sub-question
sketch (\texttt{dynamical\_substrate\_law/sketches/F1\_subquestion\_pcd\_orientation\_link.md})
has developed it. We lift the two load-bearing results and recast them as a
classification statement. (Their viability-level status, and the three open commitments
that bound it, are carried forward explicitly in \S\ref{sec:layer}.)

\subsection{Phase 1: the substrate net DI-bit current is aligned with $\nhat$}

Under Mechanism A (a direction-correlated DI-bit propagation-rate asymmetry parameterised
by $\delta$ under the substrate primitive $\nhat$), the F.1 Phase-1 result registers the
substrate net DI-bit current at the host vertex (F.1 \S11.4):
\begin{equation}
\vec{j}_{DI}^{\,\text{net}}(\vhost) \;=\; \frac{6\,r_0\,\delta}{\phi^2}\,\nhat \;+\; \Oof{\delta^2}.
\label{eq:phase1}
\end{equation}
The current is \emph{aligned with $\nhat$}: its unit direction is
$\jnet(\vhost) = \mathrm{sign}(\delta)\,\nhat$, a polar, $TI$-even directional content
(\S\ref{sec:parity}). The structurally important fact for the present theorem is not the
magnitude $6r_0\delta/\phi^2$ but the \emph{direction}: the substrate furnishes no
independent direction here --- the net current points along the already-registered
primitive $\nhat$.

\subsection{Phase 2: the cycle-orientation coupling rule}

The cycle's intrinsic forward direction P$\to$C$\to$D is a $TI$-odd pseudoscalar
$\scyc \in \{+1,-1\}$ (F.1 \S12.3): $T$-odd because time-reversal reverses the cycle's
temporal sequencing, $I$-even because the cycle is internal to each CP. The F.1 Phase-2
coupling rule (F.1 \S12.5) constructs the substrate cycle-orientation pseudovector as the
product
\begin{equation}
\wPCD(v) \;:=\; \scyc \cdot \jnet(v),
\label{eq:coupling}
\end{equation}
the unique lowest-order $TI$-odd substrate object assembled from the two framework
contents (a $TI$-odd pseudoscalar times a $TI$-even unit vector). Combining
\eqref{eq:phase1}--\eqref{eq:coupling} at the host vertex with $\scyc = +1$ by global
convention and $\delta > 0$:
\begin{equation}
\boxed{\;\wPCD(\vhost) \;=\; \scyc \cdot \nhat\;}
\label{eq:result}
\end{equation}
the cycle-orientation pseudovector is the product of the cycle handedness $\scyc$ and the
substrate primitive direction $\nhat$.

\subsection{What is NOT lifted}

The F.1 Phase-3 (magnitude $|M^{\text{thermo}}| = |\delta|/6$, the Case~A.1 $\delta=\chi$
unification giving $\chi/6\approx 0.0394$) and Phase-4 (leptogenesis phenomenology at
1.64\%) are \emph{not} inputs to this theorem. They bear on the magnitude $|\chi|=\phi^{-3}$
(OPEN-CHIR-1d / audit E21) and on the SM-alignment programme (OPEN-CHIR-3), not on E20's
primitive count. Excluding them is precisely what lets the count be claimed more robustly
than the full F.1 closure (\S\ref{sec:layer}).

% ============================================================
\section{Axiom-attribution reconciliation (precondition cleared)}
\label{sec:axiom}
% ============================================================

The F.1 sketch attributes the cycle pseudoscalar $\scyc$ to ``CPP axiom A5 or
equivalent'' (F.1 \S12.3). In the canonical post-consolidation \emph{nine}-axiom set
(\texttt{axiom-registry.md}), \textbf{A5 is the metric} ($\eta = \ell_{\text{edge}}/
R_{\text{circ}} = 1/\phi$), which carries no cycle-direction content. The PCD cycle's
ordering and forward direction are content of \textbf{A1} (the CP's capacity to perceive
and respond, defining the Perceive/Compute/Displace steps) $+$ \textbf{A4} (the Nexus,
enforcing lattice-wide consistency at each Absolute Moment, sequencing the steps). This is
exactly the audit's temporal-primitive attribution (E2/E5 $=$ A1$+$A4; E17 $=$ D1 PCD
arrow). The F.1 ``A5'' is a pre-canonical-numbering attribution drift, of the same class
as the axiom-count and PCD-terminology drifts the audit cleared before shipping.

\begin{remark}[Reconciliation]
This artifact attributes $\scyc$ to \textbf{A1$+$A4} (the temporal primitive). Under this
attribution, \eqref{eq:result} reads: the cycle-orientation pseudovector is the product of
the temporal primitive's handedness (A1$+$A4) and the spatial primitive direction (E16).
Both factors are registered primitives; neither is A5. The F.1 sketch is a prior immutable
workstream and is not edited by this artifact; the reconciliation is recorded here and in
the Patch-0635 scope sketch \S5.1.
\end{remark}

% ============================================================
\section{Factor classification}
\label{sec:factors}
% ============================================================

The decomposition \eqref{eq:result} has two factors. We classify each against the
registered primitives of \S\ref{sec:setup}.

\paragraph{Axis $\nhat$ $\to$ audit E16 (spatial primitive).} The axis of $\wPCD$ is
$\nhat$ itself, inherited from the net DI-bit current \eqref{eq:phase1} being aligned with
$\nhat$. It is \emph{not an independent direction}: no new direction in $\Reals^4$ is
introduced; the cycle orients along the already-registered substrate primitive. Any
residual freedom (the sign of $\delta$, the $\pm\nhat$ choice) is absorbed into the
sign factor below, not into a new axis.

\paragraph{Sign $\scyc$ $\to$ audit E2/E5/E17 (temporal primitive).} The sign of $\wPCD$
is the cycle handedness $\scyc$, the pseudoscalar content of the temporal primitive
(A1$+$A4, \S\ref{sec:axiom}). It is \emph{not an independent handedness}: it is the
already-registered cycle ordering expressed as a $\{+1,-1\}$-valued pseudoscalar. The
framework's commitment to a definite cycle direction (P$\to$C$\to$D by convention) is what
gives $\scyc$ its non-zero $TI$-odd content; this commitment is the temporal primitive,
not a new input.

\subsection{Two-factor $TI$-parity check}
\label{sec:parity}

The product structure \eqref{eq:result} carries the parity the manifestation-(iv)
Wigner--Eckart datum requires (F.1 \S12.6), confirming the two factors compose correctly:

\begin{center}
\begin{tabular}{@{}lccc@{}}
\hline
\textbf{Quantity} & $T$ & $I$ & $TI$ \\
\hline
$\scyc$ (temporal-primitive handedness, A1$+$A4) & $-$ & $+$ & $-$ \\
$\jnet(\vhost) = \mathrm{sign}(\delta)\,\nhat$ (spatial primitive, E16) & $-$ & $-$ & $+$ \\
\textbf{Product} $\wPCD(\vhost) = \scyc\,\nhat$ & $+$ & $-$ & $-$ \\
\hline
\end{tabular}
\end{center}

The product is $T$-even, $I$-odd, $TI$-odd. The two factors carry orthogonal framework
content --- one purely temporal ($\scyc$, the cycle handedness), one purely spatial
($\nhat$, the substrate direction) --- and their product is the cycle-orientation. No
third content is needed or present.

% ============================================================
\section{The primitive-count theorem}
\label{sec:count}
% ============================================================

\begin{theorem}[$\nhat \mapsto \wPCD$ primitive count; THEO-CHIR-PCD-ORIENTATION-1]
\label{thm:count}
Under CPP axioms A1--A11 $+$ Reading~C ($\nhat$ as substrate primitive) $+$ Mechanism~A,
the PCD-cycle orientation pseudovector at the host vertex satisfies
\[
\wPCD(\vhost) \;=\; \scyc \cdot \nhat,
\]
a product of the registered spatial primitive $\nhat$ (audit E16) and the registered
temporal primitive's handedness $\scyc$ (audit E2/E5/E17, A1$+$A4). Consequently $\wPCD$
introduces \emph{no independent third chirality primitive}: it is a function of the two
already-registered primitives alone. \textbf{Scenario~B (}$\wPCD$ \textbf{an independent
primitive parallel to} $\nhat$\textbf{) is refuted}, and audit entry E20 is \emph{emergent}
from the two registered primitives jointly. The audit's headline finding --- all spatial
chirality reduces to $\nhat$ --- is preserved, since $\wPCD$ contributes no spatial
direction beyond $\nhat$.
\end{theorem}

\begin{proof}[Proof (factorization).]
By \S\ref{sec:decomp}, $\wPCD(\vhost) = \scyc\cdot\jnet(\vhost)$ (coupling rule
\eqref{eq:coupling}) and $\jnet(\vhost) = \mathrm{sign}(\delta)\,\nhat$ (Phase-1 alignment
\eqref{eq:phase1}), giving $\wPCD(\vhost) = \scyc\,\nhat$ \eqref{eq:result}. By
\S\ref{sec:factors}, the axis factor is the registered primitive $\nhat$ (E16) and the sign
factor is the registered temporal primitive's handedness $\scyc$ (E2/E5/E17, attributed to
A1$+$A4 per \S\ref{sec:axiom}). Both factors are registered; their product introduces no
new direction (the axis is $\nhat$) and no new handedness (the sign is $\scyc$). Hence
$\wPCD$ is a function of $\{\nhat,\scyc\}$ alone, with no third independent primitive.
Scenario~B requires $\wPCD$ to be an independent primitive; the factorization exhibits it
as a function of two registered primitives, contradicting independence. Therefore Scenario~B
is refuted and E20 is emergent from $\{\nhat,\scyc\}$.
\end{proof}

\subsection{Robustness to the F.1 open commitments}
\label{sec:robust}

Theorem~\ref{thm:count} is a primitive-\emph{count} statement. The F.1 trajectory's three
open structural commitments (F.1 \S14.17) bound the \emph{strength} of the underlying
substrate result but, crucially, none of them reintroduces an independent primitive:

\begin{center}
\begin{tabular}{@{}p{4.6cm}p{3.0cm}p{6.0cm}@{}}
\hline
\textbf{F.1 open commitment} & \textbf{What it affects} & \textbf{Effect on the primitive count} \\
\hline
Local algebraic representation ansatz (Phase 2) & the \emph{representation} of the observable & none --- a different representation still yields $\wPCD$ as a function of $\nhat$ and the cycle handedness; no new direction/handedness \\
$\scyc$ algebraization (explicitation of latent cycle structure) & whether $\scyc$ is \emph{explicitated} vs latent in the temporal primitive & none --- in either case $\scyc$ is the temporal primitive (A1$+$A4), not a new primitive \\
Case~A.1 ($\delta=\chi$) (Phase 3) & the \emph{magnitude} $|\chi|=\phi^{-3}$ & none on E20 --- magnitude is audit E21 / OPEN-CHIR-1d, a separate entry \\
\hline
\end{tabular}
\end{center}

Even under the worst case for all three commitments, $\wPCD$ remains a function of
$\{\nhat,\scyc\}$ and introduces no third primitive. The primitive count is therefore firm
where the magnitude and Layer-3 rigor are provisional. This is the load-bearing economy of
resolving E20 as a classification (primitive-count) theorem rather than as the full F.1
magnitude closure.

% ============================================================
\section{Layer, reclassification, exclusions, and falsifiers}
\label{sec:layer}
% ============================================================

\subsection{Layer: viability (2.5), inheriting the F.1 ceiling}

The decomposition \eqref{eq:result} is lifted from the F.1 Phase-1/Phase-2 results, which
the F.1 trajectory itself rates at \emph{viability level} (Layer~2.5), not Layer~3 (F.1
\S14.17). This theorem does not exceed that ceiling. The primitive-count conclusion
(Scenario~B refuted; $\wPCD$ a function of two registered primitives) is robust to the F.1
open commitments (\S\ref{sec:robust}), but the underlying decomposition is viability-level,
so:

\begin{center}
\fbox{\parbox{0.92\textwidth}{\centering
\textbf{Reclassification of audit E20:} \emph{unregistered (conditional)}
$\;\longrightarrow\;$ \textbf{emergent (provisional)}.\\[2pt]
Emergent, because $\wPCD$ reduces to the two registered primitives (no independent
primitive). Provisional, not established, because the reduction rests on the F.1
viability-level result, not a Layer-3 derivation.}}
\end{center}

This is exactly the v1.1 emergent (E)/(P) grading instrument: E20 moves out of
\emph{unregistered} (there is now a registered reduction with no second primitive
apparent) into emergent~(P) (the reduction's Layer-3 rigor is owed). It does \emph{not}
reach emergent~(E); claiming so would be the overclaim the audit's multi-AI review cycle
corrected for the other provisional rows.

\subsection{What this theorem does not claim}

\begin{description}
\item[D1.] It does not close the full F.1 Scenario-A trajectory; the F.1 status remains
PROVISIONAL CLOSURE at viability level, with its three open commitments untouched here.
\item[D2.] It does not derive the magnitude $|\chi| = \phi^{-3}$ (audit E21 /
OPEN-CHIR-1d) or any observable value; \S\ref{sec:decomp} excludes the F.1 magnitude content.
\item[D3.] It does not derive $\nhat$ or $\scyc$; both remain primitives. The theorem shows
$\wPCD$ is their product, not that either factor is itself reducible.
\item[D4.] It does not resolve the capture handedness E19 (\S\ref{sec:e19}).
\end{description}

\subsection{Cross-link to E19, flagged not resolved}
\label{sec:e19}

The sign factor $\scyc$ is the \emph{temporal} cycle handedness. It must be kept distinct
from the audit's deepest unregistered entry E19 (D3 capture handedness), which is a
\emph{spatial/CP} handedness in the dipole-pairing rule, consumed by SD-CHIR-1/2 via
$\zeta^W$ and $\zeta^{qDP}$. If a future result identifies $\scyc$ with the capture
handedness, that would \emph{merge} E20 into E19 (reducing two items to one --- a
programme-positive reclassification), not refute the present theorem. The merge is flagged
as an open cross-link (and as falsifier F3 below in its merge direction), not asserted.

\subsection{Exclusions and anti-erasure}

\begin{description}
\item[X1.] The theorem decides the primitive count of E20 only; it does not address the
magnitude (E21), the SM-alignment chain (E26 / OPEN-CHIR-3), or manifestation~(v).
\item[X2.] Mechanism dependence: the axis-alignment \eqref{eq:phase1} rests on Mechanism~A
(the F.1 leading candidate). Under Mechanisms~A or~B the cycle-orientation aligns with
$\nhat$, so the primitive-count conclusion is unchanged; the theorem does not hinge on the
A-vs-B selection. Mechanism~C (clock-skew) was registered by F.1 as unmotivated and is not
relied upon.
\item[Anti-erasure.] The F.1 ``$\scyc$ from A5'' attribution is preserved for the record as
a pre-canonical-numbering drift; this artifact uses the canonical A1$+$A4 attribution
(\S\ref{sec:axiom}). An external reader encountering the legacy ``A5'' form in the F.1
sketch should consult this reconciliation.
\end{description}

\subsection{Falsifiers}

\textbf{(F1)} Exhibit an independent direction in $\wPCD$ --- a component not equal to
$\pm\nhat$ --- demonstrating an independent spatial primitive (which would restore part of
Scenario~B and break the audit's spatial-reduction headline). \textbf{(F2)} Exhibit a sign
content of $\wPCD$ independent of the temporal primitive's ordering --- a handedness not
reducible to $\scyc$ (A1$+$A4) --- demonstrating an independent temporal/CP primitive.
\textbf{(F3)} Show $\scyc$ is the capture handedness E19 in disguise; this does not refute
the theorem but \emph{merges} E20 into E19 (\S\ref{sec:e19}), a reclassification reducing
two chirality items to one.

% ============================================================
\section*{References (internal)}
\addcontentsline{toc}{section}{References (internal)}
% ============================================================

\begin{itemize}[leftmargin=1.4em]
\item \texttt{chirality\_audit/theo\_chir\_audit\_1.tex} (v1.1) --- audit entry E20 / D4
(the conditional cell); the registered primitives E16 ($\nhat$) and E2/E5/E17 (temporal);
exclusion X2; the v1.1 emergent (E)/(P) grading.
\item \texttt{chirality\_derivations/sketches/theo\_chir\_pcd\_orientation\_1\_scope.md}
--- the Patch 0635 scope-and-precondition sketch (this artifact's blueprint), \S5.1
axiom reconciliation, \S5.3 E19 cross-link.
\item \texttt{dynamical\_substrate\_law/sketches/F1\_subquestion\_pcd\_orientation\_link.md}
--- \S11 Phase~1 (net DI-bit current $\parallel\nhat$, DSL-1$\to$DSL-4); \S12 Phase~2
(coupling rule $\wPCD = \scyc\,\jnet$); \S6 (Scenario A/B/C); \S14.17 (PROVISIONAL CLOSURE
at viability level --- the inherited ceiling).
\item \texttt{axiom-registry.md} --- canonical nine axioms; A1 (CP perceive/respond) $+$
A4 (Nexus / Absolute-Moment cadence) carry the cycle ordering; A5 is the metric (the
\S\ref{sec:axiom} reconciliation point).
\item \texttt{capotauro/} --- Reading~C ($\nhat = \vhost$, $\Hfour \to \Hthree = \Icoh$);
FI-C-RC-1; $|\chi| = \phi^{-3}$.
\item \texttt{frontier\_sectors/CHIR.md} --- OPEN-CHIR-F1-LINK; OPEN-CHIR-1d (E21
magnitude); OPEN-CHIR-1c/2d (E19 capture handedness).
\end{itemize}

\vfill
\noindent\rule{\textwidth}{0.4pt}\\
{\footnotesize THEO-CHIR-PCD-ORIENTATION-1 v1.0 --- Patch 0636, Session 148 --- Conscious
Point Physics, Substrate Chirality Arc. Viability-level (Layer~2.5) primitive-count
classification theorem resolving audit entry E20 / OPEN-CHIR-F1-LINK: $\wPCD = \scyc\,\nhat$,
a product of the registered spatial ($\nhat$, E16) and temporal ($\scyc$, A1$+$A4,
E2/E5/E17) primitives; Scenario~B refuted; E20 $\to$ emergent (provisional). Built from the
Patch 0635 scope sketch; axiom-attribution precondition (A5$\to$A1$+$A4) cleared at \S3.}

% BEGIN GENERATED IDENTIFIER APPENDIX -- do not edit by hand
\clearpage
\section*{Appendix: Programme identifiers used in this paper}
\addcontentsline{toc}{section}{Appendix: Programme identifiers}

\noindent This programme tracks its unresolved questions under short identifiers so that a claim and its open dependencies can be cited precisely. The identifiers appearing in this paper are glossed below; no external document is needed to read them.

\begin{description}
  \item[\texttt{OPEN-CHIR-1}] \hfill \\ Derive each entry the audit classifies as *emergent* from
  \item[\texttt{OPEN-CHIR-2}] \hfill \\ For each entry the audit classifies as *unregistered*,
  \item[\texttt{OPEN-CHIR-3}] \hfill \\ Connect substrate chirality (FI-C-9, `sign(δ)`) to electroweak chiral structure (parity violation, δ\_CP). Decide whether substrate chirality is the **source** of SM chiral structure (→ chirality fully emergent V2/W2 via the SM) or independent of it (→ two distinct primitives).
  \item[\texttt{OPEN-CHIR-F1-LINK}] \hfill \\ Whether the PCD winding is derived from the chirality director (Scenario A) or is an independent primitive (Scenario B).
\end{description}
% END GENERATED IDENTIFIER APPENDIX

\end{document}
